\subsection{Weak-interface trace projection: rotated bi-material bar}
\label{sec:weakTraceRotatedBar}

This folder verifies the local-stress behavior of the prescribed weak-interface
trace projection.  A diagonal two-material coupon is loaded in plane-strain
tension by moving the global x and y F-table faces.  The stiff phase is assigned
a low Von Mises yield stress, while the prescribed strain is chosen so the ideal
series response remains elastic.  A single-field mixed-cell stress error can
therefore be detected as local stiff-side stress elevation and spurious
plasticity near the inclined interface.

The active variants are: single-field baseline, false-elastic cohesive-zone
coupling, and prescribed-surface trace projection between contact groups.  DFG
surface-flag and stress-residual partitioning variants are included as
placeholders for future development.

\paragraph{Geometry convention.}  The two material segments are standard PFW
polygon objects.  They extend into the lower-left and upper-right loading
corners so the moving boundaries load finite boundary segments.  Only the
internal material contact face is assigned a nonzero surface flag in the CZ and
weak-discontinuity trace variants.

\paragraph{Diagnostics.}  The quantitative metrics use tracer particles placed
on near-interface and far-field layers in both materials.  The post-processor
computes the material-normal stress and von Mises stress from tracer stress
components.  Optional VisIt frames render both particle material regions using a
case-specific multi-region plot script.

\input{\CaseOutputDir/weak_trace_rotated_bar_results.tex}
